\section{Background}
\label{sec:background}

\subsection{NSGA-II}

\nsga{} (Non-dominated Sorting Genetic Algorithm II) was introduced by
Deb, Pratap, Agarwal, and Meyarivan~\citeyear{deb2002} as a fast, elitist
multi-objective evolutionary algorithm.
It addresses the two central challenges of multi-objective optimisation:
(1) converging to the Pareto front, and (2) maintaining a diverse spread of
solutions along the front rather than collapsing to a single region.

\nsga{} achieves both via a two-level selection criterion.
Plans are first ranked by Pareto dominance: a plan on the Pareto front
(dominated by no other plan) receives rank 0; a plan dominated only by
rank-0 plans receives rank 1; and so on.
Among plans at the same Pareto rank, \nsga{} prefers those with higher
\emph{crowding distance} --- a measure of isolation in objective space.
Crowding distance rewards plans that occupy sparsely populated regions of the
objective space, ensuring diversity along the front.

The algorithm is \emph{elitist} in the sense that the best plans from the
current generation always survive to the next (the offspring pool of size
$N_{\mathrm{pop}}$ is added to the current population, giving $2N_{\mathrm{pop}}$
candidates, and the best $N_{\mathrm{pop}}$ are retained by
rank-then-crowding-distance).
Deb et al.~\citeyear{deb2002} show empirically that this elitism substantially
improves convergence speed on standard benchmark functions.

Zitzler, Deb, and Thiele~\citeyear{zitzler2000} compare \nsga{} with six
other multi-objective algorithms on a suite of nine benchmark problems.
\nsga{} consistently achieves strong coverage of the Pareto front and is
competitive with or superior to SPEA, PAES, and MOGA on all benchmark
functions.
For problems with 2--3 objectives and moderate population sizes
($N_{\mathrm{pop}} \leq 200$), \nsga{} is the standard reference algorithm
in the multi-objective evolutionary computation literature.

\subsection{Prior Work on Multi-Objective Redistricting}

Redistricting has been treated as a single-objective optimisation problem in
the vast majority of prior work: algorithms minimise a weighted aggregate of
criteria and return a single plan.

Gurnee and Shmoys~\citeyear{gurnee2021} introduce Fairmandering, which uses
column generation to optimise a fairness-weighted sum of compactness and
partisan deviation.
Their formulation requires the practitioner to specify weights, and different
weight choices produce substantially different plans.
They do not explore the trade-off structure between their objectives --- the
Pareto front.

DeFord, Duchin, and Solomon~\citeyear{deford2021} introduce ReCom, a Markov
chain over redistricting plans that targets a near-uniform distribution over
valid plans weighted by compactness.
ReCom ensemble analysis characterises the partisan and compactness distributions
over valid plans, but does not directly optimise for multiple objectives.
The SMC-Pareto approach in this paper uses the ReCom-derived ensemble
(produced by \smc{}~\citep{dellaLibera2026smc}) as a starting point for
Pareto front construction.

The closest prior work to this paper is the concurrent stream on
evolutionary redistricting, which applies genetic algorithms to weighted-sum
redistricting objectives.
To our knowledge, this paper is the first to apply NSGA-II or any
multi-objective evolutionary algorithm to redistricting, and the first to
use the Pareto front as a legal tool via the enacted-plan dominance test.

\subsection{The Three Objectives}

Every redistricting plan is evaluated on three objectives simultaneously.
All three are minimised.

\subsubsection{EC: Edge Cuts (Compactness)}

The edge-cut objective counts the number of edges in the tract adjacency
graph that cross district boundaries.
Lower edge cut means more compact districts: fewer tract-to-tract boundaries
are split.

\begin{definition}[Edge cut]
\label{def:ec}
Let $G = (V, E)$ be the tract adjacency graph and $\pi : V \to \{1,\ldots,k\}$
a redistricting plan (assignment of tracts to districts).
The edge cut of $\pi$ is:
\[
  \EC(\pi, G) \;=\; \bigl|\{(u,v) \in E : \pi(u) \neq \pi(v)\}\bigr|.
\]
\end{definition}

Edge cut is the compactness metric used throughout this platform (B.2, B.7,
B.24, G.15).
It is the natural graph-theoretic compactness measure and is computationally
tractable as an optimisation objective: evaluating $\EC(\pi)$ for any plan
$\pi$ takes $O(|E|)$ time.

\subsubsection{$\Dseats$: Partisan Distance (Proportionality)}

The partisan distance objective measures how far the plan is from proportional
representation, using 2020 presidential precinct returns as the vote signal.
Lower $\Dseats$ means closer to proportional.

\begin{definition}[Partisan distance]
\label{def:dseats}
Let $V_D$ be the total Democratic presidential votes and $V_{\mathrm{tot}}$
the total votes cast in the state.
The proportional Democratic seat share is
$s_{\mathrm{prop}} = k \cdot (V_D / V_{\mathrm{tot}})$.
Let $D_w(\pi)$ be the number of districts won by the Democratic candidate
under plan $\pi$.
Then:
\[
  \Dseats(\pi) \;=\; |D_w(\pi) - s_{\mathrm{prop}}|.
\]
\end{definition}

The absolute-value framing is neutral: it does not favour Democratic or
Republican outcomes, only representational fidelity.
A plan that gives Democrats 3 more seats than proportional scores the same as
one that gives Republicans 3 more seats.
We use 2020 presidential two-party returns as the partisan signal; presidential returns are standard in ensemble redistricting research \citep{deford2021} but may differ from Senate or gubernatorial returns in polarised cycles --- practitioners should verify sensitivity to vote vintage.

\paragraph{Discrete resolution.}
$D_w(\pi) \in \{0, 1, \ldots, k\}$ is integer-valued.
For NC with $k=14$, $s_{\mathrm{prop}} \approx 7.2$, so
$\Dseats \in \{0.2, 0.8, 1.2, \ldots\}$ --- roughly 14 distinct values.
This discrete resolution means the Pareto ordering on $\Dseats$ alone
provides limited differentiation; the effective Pareto front may collapse
to a two-dimensional $\EC$--$\VRA$ frontier when many plans share the same
$D_w$ value.
This is disclosed in the NDJSON output via \texttt{d\_seats\_discrete: true}.
Phase~2 will introduce vote margin (continuous) as a finer-grained proxy.

\subsubsection{$\VRA$: VRA Deficit (Minority Representation)}

The VRA deficit objective measures the degree to which minority-opportunity
districts are underserved.
Lower $\VRA$ means better compliance with the Voting Rights Act.

\begin{definition}[VRA deficit]
\label{def:vra}
Let $\mu_d(\pi)$ be the minority voting-age population (VAP) share in district
$d$ under plan $\pi$.
The VRA deficit is:
\[
  \VRA(\pi) \;=\; \sum_{d=1}^{k} \max\bigl(0,\; 0.50 - \mu_d(\pi)\bigr).
\]
Only districts below the 50\% minority VAP threshold contribute.
Districts already at or above 50\% contribute zero.
\end{definition}

The 50\% threshold corresponds to the standard majority-minority district
criterion.
$\VRA = 0$ means every district that was formerly a majority-minority district
(or should be, given the minority population distribution) achieves at least
50\% minority VAP.
Practitioners may filter the NDJSON output by \texttt{vra\_deficit = 0} to
consider only VRA-compliant plans from the frontier.

\textit{Note}: this implementation counts \emph{all} districts below the 50\% minority VAP threshold, not solely legally designated majority-minority districts. States with few minority-opportunity districts may see VRA\_deficit = 0 under this definition regardless of demographic distribution. Practitioners in covered jurisdictions should apply the Gingles feasibility analysis to determine which districts are legally subject to VRA Section 2 scrutiny.

\paragraph{VRA as objective, not constraint.}
The spec treats $\VRA$ as an objective rather than a hard constraint, so
that the Pareto front can display the compactness cost of VRA compliance:
plans with $\VRA > 0$ remain in the frontier if they are non-dominated, and
practitioners can observe exactly how much compactness must be sacrificed to
achieve $\VRA = 0$.
A \texttt{--vra-constrained} flag (Phase~2) will filter to VRA-compliant plans
before running NSGA-II.

\subsection{Pareto Dominance and the Pareto Front}

\begin{definition}[Pareto rank]
\label{def:paretorank}
The Pareto rank of plan $\pi$ in population $P$ is the index of the smallest
front $\mathcal{F}_i$ to which $\pi$ belongs:
\begin{align*}
  \mathcal{F}_0 &= \{\pi \in P : \nexists\, \pi' \in P \text{ s.t.\ } \pi' \succ \pi\}, \\
  \mathcal{F}_{i+1} &= \{\pi \in P \setminus \bigcup_{j \leq i} \mathcal{F}_j :
    \nexists\, \pi' \in P \setminus \bigcup_{j \leq i} \mathcal{F}_j
    \text{ s.t.\ } \pi' \succ \pi\}.
\end{align*}
Plans in $\mathcal{F}_0$ have Pareto rank 0 (the Pareto front of $P$).
\end{definition}

The algorithm returns only the plans in $\mathcal{F}_0$ of the final
population; these are mutually non-dominating and represent the
best-achievable trade-offs among the three objectives within the
population's resolution.

\begin{remark}[Population-limited Pareto front]
\label{rem:population-limited}
The returned $\mathcal{F}_0$ is the Pareto front of the final \nsga{}
population, not the true Pareto front of the feasible redistricting space.
The true Pareto front may contain plans not represented in a population of
100.
This is analogous to the relationship between the best plan found by
$\be{}(1000)$ and the global edge-cut optimum: the algorithm finds a
high-quality subset of the objective space but cannot certify completeness.
The enacted-plan dominance test makes a validity claim (the returned plans
are mutually non-dominating and certified to have the stated objective values)
rather than a completeness claim (the full Pareto front has been found).
See Section~\ref{sec:legal} for the validity vs.\ completeness distinction.
\end{remark}
